% !TEX program = pdflatex
% Example: IEEE journal article template
% Compile with: pdflatex main.tex && bibtex main && pdflatex main.tex && pdflatex main.tex
% Requires: IEEEtran class (included in TeX Live / MiKTeX)

\documentclass[journal]{IEEEtran}

\usepackage{cite}
\usepackage{amsmath,amssymb}
\usepackage{graphicx}
\usepackage{url}

\begin{document}

\title{A Novel Approach to Signal Processing Using Deep Learning}
\author{Author~Name, \IEEEmembership{Member,~IEEE}}

\maketitle

\begin{abstract}
This paper proposes a novel approach to signal processing
using deep learning architectures. We demonstrate that
convolutional neural networks can effectively extract
features from raw signals without manual preprocessing.
\end{abstract}

\begin{IEEEkeywords}
Signal processing, deep learning, neural networks, feature extraction.
\end{IEEEkeywords}

\section{Introduction}
\IEEEPARstart{S}{ignal} processing has traditionally relied
on hand-crafted features and domain expertise. Recent advances
in deep learning have shown that neural networks can learn
relevant features directly from raw data~\cite{ref1}.

\section{Methodology}
We propose a multi-layer architecture consisting of
convolutional layers followed by fully connected layers.
The input signal $x(t)$ is processed through the network
to produce an output $y = f(Wx + b)$.

\section{Results}
Table~\ref{tab:results} shows the performance comparison.

\begin{table}[htbp]
\caption{Performance comparison.}
\label{tab:results}
\centering
\begin{tabular}{lcc}
\toprule
Method & Accuracy & F1-Score \\
\midrule
Baseline & 85.2\% & 0.84 \\
Proposed & 93.7\% & 0.93 \\
\bottomrule
\end{tabular}
\end{table}

\section{Conclusion}
We presented a novel deep learning approach that outperforms
traditional methods by 8.5\% in accuracy.

\bibliographystyle{IEEEtran}
\bibliography{references}

\end{document}
